\chapter*{GENERAL INTRODUCTION}
\markboth{\MakeUppercase{GENERAL INTRODUCTION}}{}
\addcontentsline{toc}{chapter}{GENERAL INTRODUCTION}
\adjustmtc
\thispagestyle{MyStyle}

In industrial manufacturing, a Request for Quotation (RFQ) constitutes a complex, multidisciplinary engineering and operational process. It serves as the entry point of a prolonged lifecycle that requires the rigorous coordination of engineering, production, quality assurance, and procurement departments. The outcome of this process determines whether a project is acquired under viable, profitable terms or results in contractual vulnerability. The present work addresses not the arithmetic generation of the quotation itself, but the structural resilience of the surrounding process: namely coordination, operational visibility, capitalization on historical data, and the early identification of feasibility risks. These elements distinguish a systematic estimation framework from ad-hoc, individual-driven workflows lacking systemic support.

This graduation project was carried out within \hostCompany, an engineering consultancy specializing in cloud architecture, microservices, artificial intelligence, and data engineering. The study is grounded in a comprehensive on-site audit conducted at \clientCompany, a major industrial conglomerate in Dammam, Saudi Arabia, and its subsidiary GHI (Gulf Heavy Industries Co.), a leading heavy steel fabrication enterprise in the Middle East. The audit identified twelve critical operational constraints, categorized into five structural domains: process coordination, visibility, capitalization of historical data, feasibility and supplier risk management, and tooling resilience.

The initial mandate consisted of developing an artificial intelligence system for the automated generation of Bills of Quantities (BoQ). However, the audit revealed significant challenges: the tacit expertise embedded in the estimation workbooks was not readily extractable as structured training data, and the asymmetric cost of error inherent to industrial quotation precluded the deployment of a fully automated pricing system without intolerable financial risk. Consequently, the project was strategically reoriented to design a decision-support platform that fortifies all process stages independent of irreducible human expertise, while integrating with manual expert validation.

The resulting platform, designated as \textit{Itqan} (denoting mastery and precision), is a comprehensive web-based system. It manages the RFQ lifecycle, centralizes process artifacts, provides analytical dashboards, deterministically parses engineering documents, and features a constrained conversational AI copilot grounded strictly in authorized platform evidence. The system architecture is predicated on four foundational pillars: Workflow and Tracking, Communication Automation, Business Intelligence, and Predictive Analytics. It is augmented by a cross-cutting conversational layer, the \textit{RFQ Copilot}. The primary engineering contribution lies in the trust-boundary architecture, which strictly governs the copilot's behavior in response to the asymmetric cost of error in the estimation domain. Within the scope of this project, the operational workflow-tracking backbone of the first pillar is implemented, while the second and third are partially addressed. The fourth pillar is established at the foundational level via deterministic document parsing, with predictive modeling deferred to future iterations.

The report details the project across six chapters. Chapter~\ref{chap:context} presents the organizational context, the industrial framework, the audit findings, and the resulting strategic reorientation. Chapter~\ref{chap:soa} situates the proposed platform within the current technological and academic state of the art. Chapter~\ref{chap:requirements} formalizes the system actors, functional and non-functional requirements, and the development methodology. Chapter~\ref{chap:architecture} details the conceptual and technical architecture, highlighting the trust-boundary architecture as the central innovation. Chapter~\ref{chap:implementation} discusses the implementation of the platform components: the manager service, the copilot service, the intelligence service, and the user interface. Finally, Chapter~\ref{chap:validation} evaluates the system and verifies, within the documented evidence scope, that the architectural invariants are preserved under empirical evaluation. The General Conclusion summarizes the contributions and outlines perspectives for future work.
